\chapter*{Загальні висновки}
\addcontentsline{toc}{chapter}{Загальні висновки}

\vspace{-0.5em} % Змінюйте це число, щоб регулювати відступ (чим більше число, тим ближче текст)

\begin{center}
    \textit{Безпека — це межі платформи}
\end{center}

\vspace{0.01cm}

Аналіз схем шахрайства на торговельних майданчиках доводить, що сьогодні  \textbf{головним вектором атаки є експлуатація людського фактору}, а не лише технічна недосконалість платформ. Безпека онлайн-торгівлі — це не відмова від довіри, а \textbf{перехід до усвідомленої обережності}. Знання маніпулятивних схем дозволяє вам зберігати контакт із цифровим світом, не стаючи об'єктом професійного впливу.

Ваша безпека базується на непорушному правилі «цифрового периметра»: усі операції — від обговорення деталей до оплати та оформлення доставки — мають відбуватися виключно всередині офіційного додатка чи сайту. Будь-яка пропозиція перейти у сторонній месенджер є першим сигналом небезпеки, адже саме там шахраї виводять вас з-під захисту системи безпеки платформи.

Також варто переглянути своє ставлення до будь-яких форм авансових платежів у світі, де існують сервіси безпечних угод та післяплати. Пам’ятайте, що емоційний тиск, штучний ажіотаж та підозріло низька ціна — це не ознаки «вигідної пропозиції», а інструменти маніпуляції. Вони мають стати для вас автоматичним сигналом для ввімкнення критичного мислення та аналізу ситуації, що дозволить вам приймати зважені рішення без стороннього впливу.


\begin{quote}
    \itshape
    \textbf{Запам'ятайте!} Ви не залишились сам на сам із проблемою. Автор готовий допомогти у будь-якому питанні, надати фахову консультацію або порадити перевірених спеціалістів.
\end{quote}

